\section{Shared-parameter coordinate-ring decomposition}
\label{sec:coordinate-ring}

We now retain the polynomial equalities, including square-slack equations, and
separate the one shared parameter algebra from sample-local algebras.  Let
\(k=\R\) for the geometric statements and put
\begin{equation}
  A=k[\theta_1,\ldots,\theta_P].
\end{equation}
For sample \(i\), let \(S^{(i)}\) denote a finite tuple of local indeterminates
and define
\begin{equation}
  B_i=A[S^{(i)}],\qquad
  \cR^{(i)}=B_i/I^{(i)},
\end{equation}
where \(I^{(i)}\) is generated by the polynomial equalities belonging to that
sample.  In the common ambient polynomial algebra
\begin{equation}
  R=A[S^{(1)},\ldots,S^{(N)}],
\end{equation}
each \(I^{(i)}\) is understood through extension from \(B_i\) to \(R\), and
\begin{equation}
  I=\sum_{i=1}^{N} I^{(i)}R,
  \qquad
  \cR=R/I.
  \label{eq:global-ring}
\end{equation}

\begin{proposition}[Ambient shared-base tensor product]
\label{prop:ambient-tensor}
There is a canonical \(A\)-algebra isomorphism
\begin{equation}
  A[S^{(1)},\ldots,S^{(N)}]
  \cong
  A[S^{(1)}]\otimes_A\cdots\otimes_A A[S^{(N)}].
  \label{eq:ambient-tensor}
\end{equation}
\end{proposition}

\begin{proof}
Both sides are the coproduct in commutative \(A\)-algebras of the polynomial
\(A\)-algebras on the disjoint variable sets \(S^{(i)}\).  Concretely, the map
sends each variable to its copy in the corresponding tensor factor.  The
inverse multiplies the images of pure tensors in the common polynomial ring.
The two maps agree on \(A\) and on every variable, and hence are inverse by the
universal property of polynomial algebras.
\end{proof}

\begin{theorem}[Quotient tensor decomposition]
\label{thm:quotient-tensor}
With the notation above, there is a canonical \(A\)-algebra isomorphism
\begin{equation}
  \boxed{
  \cR\cong
  \cR^{(1)}\otimes_A\cdots\otimes_A\cR^{(N)}.}
  \label{eq:quotient-tensor}
\end{equation}
\end{theorem}

\begin{proof}
It suffices to treat two factors and iterate.  The canonical bilinear map
\(B_1\times B_2\to (B_1\otimes_A B_2)/J\), where \(J\) is generated by the
images of \(I^{(1)}\otimes B_2\) and \(B_1\otimes I^{(2)}\), vanishes whenever
either input belongs to the corresponding ideal.  It therefore descends to a
map
\begin{equation}
  (B_1/I^{(1)})\otimes_A(B_2/I^{(2)})
  \longrightarrow (B_1\otimes_A B_2)/J.
\end{equation}
The universal property gives an inverse, so the map is an isomorphism.  Under
\cref{eq:ambient-tensor}, the ideal \(J\) corresponds to the sum of the two
extended sample ideals.  Iteration gives \cref{eq:quotient-tensor}.
\end{proof}

\begin{remark}[One parameter algebra]
The tensor product is over \(A\).  It identifies the parameter coordinates in
all local factors.  A tensor product over \(k\) would create independent copies
of the parameters and would model a different feasibility problem.
\end{remark}

\subsection{Equality rings versus feasible real sets}

\Cref{thm:quotient-tensor} concerns polynomial equalities.  The exact decision
problem concerns real points together with the existential interpretation of
slack variables.  Passing to a quotient ring does not perform existential
projection, choose a real connected component, or decide whether real points
exist.  These distinctions matter because two systems can have closely related
differential modules while having different real feasibility.

The global real feasible set may be written
\begin{equation}
  X_I(\R)=\left\{(\theta,S^{(1)},\ldots,S^{(N)}):
  F^{(i)}(\theta,S^{(i)})=0\text{ for all }i\right\},
\end{equation}
after square-slack polynomialization.  The accepted parameter set is its
projection to \(\theta\), equivalently the intersection in
\cref{eq:projection-intersection}.  The coordinate-ring decomposition explains
the sample incidence pattern, but it does not compute this projection.
